\documentclass[12pt]{article}
\usepackage[T1]{fontenc}
\usepackage[utf8]{inputenc}
\usepackage{lmodern}
\usepackage{textcomp}
\usepackage{enumitem}
\usepackage{geometry}
\geometry{margin=1in}
\begin{document}
\begin{center}
Answer Key - Sociology - Undergraduate Level Assessment 1 \
\small (Answers are ordered exactly as in the question paper.)
\end{center}

\section*{Multiple Choice}
\begin{enumerate}[label=\arabic*.]
\item \textbf{Answer:} C
\\ \textit{Options:} A) A primary group; B) A secondary group; C) An aggregate; D) A social category
\\ \textit{Note:} The correct answer is "an aggregate" because an aggregate refers to a collection of individuals who are physically present in the same place at the same time but do not interact or share a common identity or purpose, as exemplified by the seven people waiting silently for a bus.
\item \textbf{Answer:} D
\\ \textit{Options:} A) formed an inferior 'race' with low levels of intelligence; B) lived morally unsound lives of crime and squalor; C) were too reliant upon welfare benefits; D) all of the above
\\ \textit{Note:} Murray believed that the 'underclass' was characterized by a combination of factors including economic disadvantage, social dislocation, and cultural issues, which collectively define the challenges faced by this group.
\item \textbf{Answer:} B
\\ \textit{Options:} A) repress it, by promoting only the interests of elite groups; B) revive it, by reaffirming a commitment to freedom of speech; C) reproduce it, by emphasizing face-to-face contact with peer groups; D) replace it with a superior form of communication
\\ \textit{Note:} The correct answer highlights that the Internet has revitalized the public sphere by providing a platform for diverse voices and promoting freedom of speech, which fosters democratic discourse and civic engagement.
\item \textbf{Answer:} D
\\ \textit{Options:} A) the ruling class, or bourgeoisie, who exploit the proletariat; B) a capitalist class dependent on property ownership and advantaged life chances; C) an alignment of classes with shared interests but no state power; D) a state elite whose members are drawn overwhelmingly from a power bloc
\\ \textit{Note:} The term 'power elite' as introduced by Scott (1991) accurately describes a small group of individuals within a state who hold significant power and influence, primarily originating from a dominant power bloc, thereby reinforcing the notion of concentrated authority in sociopolitical structures.
\item \textbf{Answer:} C
\\ \textit{Options:} A) disappointment and disproportion; B) disbelief and disintegration; C) disengagement and disenchantment; D) distribution and distillation
\\ \textit{Note:} Secularization refers to the process by which religious institutions, practices, and beliefs lose their social significance, which is captured by the concepts of disengagement from traditional religious authority and disenchantment with supernatural explanations, reflecting a shift towards rational and empirical understandings of the world.
\end{enumerate}

\section*{True / False}
\begin{enumerate}[label=\arabic*.]
\setcounter{enumi}{5}
\item \textbf{Answer:} True
\\ \textit{Options:} A) True; B) False
\\ \textit{Note:} Stage 3 of the health transition is characterized by a shift from infectious diseases to chronic, degenerative diseases as the primary health challenges, reflecting improvements in healthcare and living conditions that extend life expectancy but lead to increased prevalence of conditions like cancer and heart disease.
\item \textbf{Answer:} False
\\ \textit{Options:} A) True; B) False
\\ \textit{Note:} Research indicates that while women generally report higher rates of depression, individuals from lower socioeconomic backgrounds do not consistently self-report higher levels of depression compared to other groups, suggesting that socioeconomic status alone does not determine self-reported depression levels.
\item \textbf{Answer:} False
\\ \textit{Options:} A) True; B) False
\\ \textit{Note:} The statement is false because the Mafia is primarily known for its involvement in organized crime, which includes violent activities such as racketeering, extortion, and drug trafficking, rather than solely engaging in white-collar financial offenses.
\item \textbf{Answer:} True
\\ \textit{Options:} A) True; B) False
\\ \textit{Note:} Mosca and Pareto both argued that a small, organized elite class consistently dominates political power and decision-making in society, thus excluding the majority from these top positions.
\item \textbf{Answer:} True
\\ \textit{Options:} A) True; B) False
\\ \textit{Note:} A document is deemed 'authentic' when it can be verified as an original or a trustworthy reproduction attributed to its legitimate creator, ensuring its credibility and integrity in sociological research and analysis.
\end{enumerate}

\section*{Long Form Response}
\begin{enumerate}[label=\arabic*.]
\setcounter{enumi}{10}
\item \textbf{Answer:} Booth's (1901) study on poverty in London revealed a staggering poverty rate of 30.7\%, which has significant implications for understanding the social conditions of the time. This figure indicates that nearly one-third of the population lived in conditions of economic deprivation, highlighting systemic issues such as inadequate wages, lack of social support, and the challenges faced by the working class. The study not only illuminated the harsh realities of life for many Londoners but also prompted discussions about social reform and the need for policy interventions. Furthermore, Booth's work laid the groundwork for future sociological research on poverty, emphasizing the importance of empirical data in shaping our understanding of societal issues. Thus, the 30.7\% poverty rate serves as a critical lens through which we can analyze the socio-economic landscape of early 20 th-century London.
\\ \textit{Note:} The answer correctly identifies the 30.7\% poverty rate as a critical indicator of widespread economic deprivation, underscoring systemic issues affecting the working class and illuminating the social challenges prevalent in London during that period.
\item \textbf{Answer:} Warner's findings on the 'colour line' in Natchez reveal a stark division between black and white communities, illustrating the rigid caste system that defined social interactions and hierarchies. This division is not merely a reflection of physical separation but is deeply rooted in the prevailing beliefs of racial superiority held by the white population. The social, economic, and political structures in Natchez were designed to maintain this hierarchy, which justified discriminatory practices and perpetuated inequality. Warner's analysis emphasizes how these entrenched beliefs not only marginalized black citizens but also reinforced the self-image of white individuals as inherently superior, thus solidifying the colour line as a fundamental aspect of social life in the region.
\\ \textit{Note:} correct because it accurately describes how Warner's findings illustrate the entrenched social divisions and hierarchical structures in Natchez, which were underpinned by the pervasive beliefs of racial superiority among the white community.
\end{enumerate}

\vfill
\noindent\textit{End of Answer Key}
\end{document}